% ===== 原文 pp.8--27 =====
% [原文 p.8]
\section{拓扑 K 理论}

拓扑 K 理论是 K 理论中历史上最先发展起来的分支，由 Atiyah 与 Hirzebruch 在 1960 年代创立。拓扑 K 理论最著名的两个应用，也许就是 Adams 关于球面上连续线性无关切向量场最大数目的结果，以及这样一个事实：$\R$ 上除维数 $1,2,4,8$ 之外不存在有限维可除代数，这四个维数分别对应实数、复数、四元数与八元数。本章的基本参考文献是 M.~F.~Atiyah 在 1960 年代所写的经典 Harvard 讲义。此后文献中出现的其他一些书与讲义都以 Atiyah 的讲义为基础。核心的概念是向量丛的概念。

% [Tutor Note: p.8 原文 "Perhaps the two most famous applications ... is the result" 主谓不一致（applications 为复数却用 is）；照译不改。另 "clasic" 应为 classic、"Mobious" 应为 Möbius，均系原文拼写笔误，译文按正确拼写处理。]

\section{第一讲（向量丛的拓扑预备）}

大家都熟悉圆环带（annulus）与 Möbius 带。这两个空间都可以由圆周与直线构造出来：前者是二者真正的笛卡尔积，而后者是它们的「扭曲」乘积。我们从向量丛的基本定义开始；这个定义把上述两种构造加以推广：把圆周换成任意拓扑空间，把直线换成某个有限维向量空间。我们用符号 $\F$ 表示实数域 $\R$ 或复数域 $\C$。

% [Tutor Note: p.8 原文 "the field of real R" 是 "the field of real numbers R" 的漏字；照译。另注意本节此后所有丛都取 $\F=\R$ 或 $\C$，而定义 1 的纤维维数 $n$ 与 $E$ 的拓扑维数可能不同，原文已作提醒。]

\begin{defn}{1}
一个\emph{向量丛} $\xi=(E,\pi,B)$ 由下列数据组成：
\begin{enumerate}[label=\arabic*.]
\item 一个拓扑空间 $E$（也记作 $E(\xi)$，通常假定为流形），称为\emph{全空间}（total space）。
\item 一个拓扑空间 $B$（通常假定为流形），称为\emph{底空间}（base space）或简称\emph{底}（base）。
\item 一个连续满射 $\pi: E\to B$，称为\emph{投影}（projection），它满足下列性质：
\begin{enumerate}[label=\roman*.]
\item 对每个 $b\in B$，空间 $E_b(\xi)=F_b(\xi)=\pi^{-1}(b)\subset E$——称为 $b$ 上的\emph{纤维}（fibre）——具有有限维向量空间的结构，维数设为
% [原文 p.9]
$n$。此时我们说该向量丛的维数（或秩）为 $n$（这里要提醒一句：这个维数可能不同于 $E$ 的拓扑维数）。
\item 对每个 $b\in B$，\emph{局部平凡性条件}（local triviality condition）成立，即：若 $U_b\subset B$ 是 $b$ 的邻域，则存在映射
\[ h_b: U_b\times\F^n\to\pi^{-1}(U_b)\subset E \]
它是同胚。像上面这样的映射 $h_b$ 称为\emph{局部平凡化}（local trivialisations），它们具有这样的性质：限制到每条纤维上时给出向量空间的线性同构，即映射
\[ \F^n\to F_b(\xi) \]
其中
\[ v\mapsto h_b(b,v), \]
都是向量空间的线性同构。
\end{enumerate}
\end{enumerate}
\end{defn}

% [Tutor Note: p.9 定义 1 的 i、ii 两条只属于第 3 条的「性质」，不是与 1、2 并列的丛数据；排版上原文 i. 与 ii. 缩进更深。核对时可看 200 dpi 页图 p.8 末行至 p.9 首段。]

选择实数或复数，我们便分别得到实向量丛或复向量丛。

本章通篇为简单起见，向量丛或者用小写希腊字母表示，或者只写出投影映射，甚至在不会引起歧义时只写出全空间（用大写拉丁字母表示）。

向量丛映射可以按显然的方式定义：若 $\pi_1: E_1\to B$ 与 $\pi_2: E_2\to B$ 是同一底空间 $B$ 上的两个向量丛，则\emph{向量丛映射} $f: E_1\to E_2$ 是一个同胚，使得 $\pi_2 f=\pi_1$，并且限制到每条纤维上时给出向量空间的 $\F$-同态。显然，$B$ 上的 $\F$ 向量丛及其映射构成一个加法范畴。我们更感兴趣的是向量丛映射中一个特殊的类，即同构；相应的向量丛则称为同构的：

\begin{defn}{2}
同一底空间 $B$ 上的两个向量丛 $\xi,\eta$ 称为\emph{同构的}，记作 $\xi\simeq\eta$，如果它们之间存在同构，即它们全空间之间的一个同胚，设为 $\phi: E(\xi)\to E(\eta)$，
% [原文 p.10]
它在每条纤维上的限制 $\phi_b: F_b(\xi)\to F_b(\eta)$ 给出向量空间之间的线性同构，其中 $F_b(\xi)$ 是丛 $\xi$ 在点 $b$ 上纤维的另一种记号。显然，同构的向量丛必有相同的维数（但反之不真）。
\end{defn}

% [Tutor Note: p.9 定义 2 里 $\phi_b$ 只是 $\phi$ 在纤维上的限制；原文把同构定义为「全空间之间的同胚且逐纤维线性同构」，这比单说 $\xi\simeq\eta$ 多出的正是保持投影（$\pi_\eta\phi=\pi_\xi$）这一条，验证时须逐纤维核对。]

\begin{rem}{1}
给定任意向量丛 $\pi: E\to B$，还有另一种用\emph{粘合函数}（gluing functions）重构它的方式：取 $B$ 的一个开覆盖 $\{U_a\}$，带有局部平凡化 $h_a: \pi^{-1}(U_a)\to U_a\times\F^n$，然后把 $E$ 重构为不交并
\[ \bigcup_a (U_a\times\F^n) \]
的商空间，其作法是把 $(x,v)\in U_a\times\F^n$ 与 $h_bh_a^{-1}(x,v)\in U_b\times\F^n$ 等同，只要 $x\in U_a\cap U_b$。函数 $h_bh_a^{-1}$ 可以看作映射 $g_{ab}: U_a\cap U_b\to GL_n(\F)$，它们满足\emph{上链条件}（cocycle condition）
\[ g_{cb}g_{ba}=g_{ca} \]
于 $U_a\cap U_b\cap U_c$ 上。任何满足上链条件的粘合函数族都可以用来构造一个向量丛 $\pi: E\to B$。
\end{rem}

% [Tutor Note: p.10 上链条件 $g_{cb}g_{ba}=g_{ca}$ 的复合顺序与 $g_{ab}=h_bh_a^{-1}$ 的定义配套：在 $U_a\cap U_b\cap U_c$ 上算得 $h_ch_b^{-1}\cdot h_bh_a^{-1}=h_ch_a^{-1}$，中间的 $h_bh_b^{-1}$ 消去，恰为 $g_{ca}$。核对时把三个平凡化都限制到三重交上即可。]

\begin{rem}{2}
上面的向量丛定义可以推广：假定底空间各点上方的模型纤维一般同胚于某个固定的拓扑空间 $F$，而不是 $n$ 维向量空间，就得到以 $F$ 为纤维的任意\emph{纤维丛}（fibre bundle with fibre $F$）的定义。这样，纤维丛给出了笛卡尔积的一种推广。
\end{rem}

\begin{rem}{3}
有些作者放宽「向量丛中所有纤维都应有相同维数」这一条件；然而由于局部平凡化连续，纤维的维数必须是局部常值的。若底空间连通（本稿中总是如此），则纤维的维数将是常数（于是我们可以定义向量丛的维数或秩）。
\end{rem}

\para{例}

% [原文 p.11]
\begin{enumerate}[label=\arabic*.]
\item 给定任意拓扑空间 $B$，$B$ 上 $n$ 维\emph{平凡向量丛}（trivial vector bundle）是这样一个向量丛：全空间为 $B\times\F^n$，底空间为 $B$，而投影 $\pi$ 是到第一个因子的投影
\[ \pi: B\times\F^n\to B. \]
对于 $B$ 上 $n$ 维平凡向量丛，我们也可以用另一种记号 $\F^n_B$。（有时若不会引起混淆，我们可以完全省略底空间以简化记号。）

从这个例子可以看出，任意向量丛不必整体平凡；但平凡性条件说的是所有向量丛局部上「看起来像」乘积丛。换言之，可以说向量丛是某个（底）空间上的一族连续向量空间。

\item 设 $I=[0,1]$ 表示单位区间，$E$ 为 $I\times\R$ 在等同关系 $(0,t)\sim(1,-t)$ 下的商空间。则投影 $I\times\R\to I$ 诱导出映射 $\pi: E\to S^1$，它是一个 1 维向量丛（这类丛特称为\emph{线丛}）。由于 $E$ 同胚于挖去边界圆周的 Möbius 带，它被称为 \emph{Möbius} 丛。

% [Tutor Note: p.11 例 2 的等同关系 $(0,t)\sim(1,-t)$ 中负号是关键：若改写成 $(0,t)\sim(1,t)$ 得到的是平凡线丛 $S^1\times\R$；负号使纤维绕行一周后定向反转，这正是 Möbius 丛非平凡的原因。]

\item 设 $M$ 是维数为 $n$ 的实可微流形。则它的\emph{切丛}（tangent bundle）
\[ TM=\{(x,v)\in M\times\R^n : v\in T_xM\}, \]
其中 $T_xM$ 表示点 $x$ 处的切空间，是底空间为 $M$、纤维同构于 $\R^n$ 的 $n$ 维实向量丛。这一构造对复流形同样适用。

\item 实射影 $n$ 空间 $\R P^n$ 按定义是 $\R^{n+1}$ 中过原点的直线构成的空间。于是 $\R P^n$ 可以看作 $S^n$（$n$ 维球面）把对径点对等同后得到的商空间。$\R P^n$ 上的\emph{万有}（或「典范」）线丛以 $\R P^n\times\R^{n+1}$ 的子集 $E$ 为全空间，$E$ 定义为
\[ E=\{(x,v)\in\R P^n\times\R^{n+1} : v\in x\}. \]
于是有一个显然的投影映射
\[ \pi: E\to\R P^n \]
% [原文 p.12]
其中 $\pi(x,v)=x$（到第一个因子的投影），即 $\pi$ 的每条纤维是欧氏 $(n+1)$ 空间 $\R^{n+1}$ 内的直线 $x$（因此这个万有线丛的维数为 1，从而是线丛）。

$\R P^n$ 上的万有线丛还有一个\emph{正交补}（orthogonal complement）向量丛，其全空间为 $E^\perp$，其中
\[ E^\perp=\{(x,v)\in\R P^n\times\R^{n+1} : v\perp x\}. \]
它有显然的到第一个因子的投影。这个向量丛的维数为 $n$。

% [Tutor Note: p.12 「维数为 $n$」可核对：纤维 $E^\perp_x=\{v\in\R^{n+1} : v\perp x\}$ 是直线 $x$ 在 $\R^{n+1}$ 中的正交补，维数 $(n+1)-1=n$，与原文一致。]

这两种构造都可以照搬到复的情形。

\item 实射影空间的自然推广是实 Grassmann 流形 $G_k(\R^n)$，它定义为 $\R^n$ 中过原点的 $k$ 维平面构成的空间（显然 $k<n$）。类似地可以定义 $G_k(\R^n)$ 上的\emph{典范} $k$ 维向量丛，它由点对 $(x,v)$ 组成，其中 $x$ 是 Grassmann 流形中的一个「点」（一个 $k$ 维子空间），$v$ 是 $x$ 中的一个向量。它也有一个正交补，是 Grassmann 流形 $G_k(\R^n)$ 上的 $(n-k)$ 维向量丛。它们也有复的版本。
\end{enumerate}

\begin{defn}{3}
给定向量丛 $\pi: E\to B$，取子空间 $A\subset B$。则 $\pi: \pi^{-1}(A)\to A$ 显然是一个向量丛，称为 $E$ 在 $A$ 上的\emph{限制}（restriction），记作 $E|_A$。
\end{defn}

\begin{defn}{4}
向量丛 $\pi: E\to B$ 的\emph{向量子丛}（vector subbundle）是一个子空间 $E_0\subset E$，它与每条纤维的交都是向量子空间，使得 $\pi: E_0\to B$ 仍是一个向量丛。
\end{defn}

\begin{defn}{5}
给定两个向量丛 $\pi_1: E_1\to B_1$ 与 $\pi_2: E_2\to B_2$，\emph{乘积向量丛}是 $\pi_1\times\pi_2: E_1\times E_2\to B_1\times B_2$，其纤维是形如 $\pi_1^{-1}(b_1)\times\pi_2^{-1}(b_2)$ 的笛卡尔积，而 $h_a\times h_b$ 是局部平凡化，其中 $h_a: U_a\times\F^n\to\pi_1^{-1}(U_a)$ 与 $h_b: U_b\times\F^m\to\pi_2^{-1}(U_b)$ 分别是 $E_1$ 与 $E_2$ 的局部平凡化。
\end{defn}

\begin{defn}{6}
给定任意向量丛 $(E,\pi,B)$，$E$ 上的一个\emph{截面}（section）是连续映射 $s: B\to E$，使得对所有 $b\in B$ 都有 $\pi(s(b))=b$。
\end{defn}

% [原文 p.13]
每个向量丛都有一个标准截面，即\emph{零截面}（zero section），它在每条纤维上取零值。人们常把零截面与它的像等同起来，后者是 $E$ 的一个子空间，经 $\pi$ 同胚地投影到 $B$ 上。有时可以通过考察零截面的补集来区分不同构的向量丛，因为任何向量丛同构 $h: E_1\to E_2$ 必把 $E_1$ 的零截面送到 $E_2$ 的零截面，从而二者的补集必同胚。与零截面相对的另一个极端是\emph{无处为零的截面}（nowhere vanishing section）。同构也把无处为零的截面送到无处为零的截面。显然平凡丛有这样的截面，但并非所有向量丛都有。例如我们知道，球面 $S^n$ 的切丛有处处非零的截面当且仅当 $n$ 为奇数。由此推出：$n$ 为偶数时 $S^n$ 的切丛不是平凡的（证明可参见 [9]）。

% [Tutor Note: p.13 「$n$ 为偶数时 $S^n$ 的切丛非平凡」的论证是：若切丛平凡则存在 $n$ 个处处线性无关的切向量场，特别地存在一个处处非零的切向量场，而这要求 $n$ 为奇数——与 $n$ 偶矛盾。原文只给了该书引用的出处 [9]，未写这一句话的反证。]

事实上，$n$ 维向量丛 $\pi: E\to B$ 同构于平凡丛，当且仅当它有 $n$ 个截面 $s_1,\dots,s_n$，使得向量 $s_1(b),\dots,s_n(b)$ 在每条纤维 $\pi^{-1}(b)$ 中线性无关。一个方向是显然的，因为平凡丛当然有截面，而同构把线性无关的截面族送到线性无关的截面族。反过来，若已有 $n$ 个线性无关的截面 $s_i$，则由
\[ h(b,t_1,t_2,\dots,t_n)=\sum_i t_i s_i(b) \]
给出的映射 $B\times\F^n\to E$ 在每条纤维上是线性同构；此外它连续，因为它与局部平凡化 $\pi^{-1}(U)\to U\times\F^n$ 的复合连续。于是 $h$ 是同构，因为它把纤维同构地映到纤维。

% [Tutor Note: p.13 「$h$ 把纤维同构地映到纤维」一句省略了逆映射连续性的论证：由 $h$ 与局部平凡化的复合连续、且逐纤维为线性同构，可在每个局部平凡化卡上把 $h$ 写成 $U\times\F^n$ 上逐纤维可逆的连续矩阵值映射，其逆亦连续。核对时可逐卡验证。]

由向量丛的定义立即得知：对任意 $b\in B$，存在 $b$ 的一个邻域 $U$ 以及 $E$ 在 $U$ 上的截面 $s_1,\dots,s_n$，使得对一切 $x\in U$，$s_1(x),\dots,s_n(x)$ 构成纤维 $F_x=\pi^{-1}(x)$ 的一组基。我们说 $s_1,\dots,s_n$ 在 $b$ 处构成\emph{局部基}（local basis），并且 $E$ 的任意截面都可写成 $s(x)=\sum_i a_i(x)s_i(x)$，其中 $a_i(x)\in\F$。显然，若所有 $a_i$ 函数连续，则 $s$ 连续。

% [Tutor Note: p.13 原文 "Clearly s is continuous of all the a_i functions are." 中 "of" 应为 "if"；此处照译作条件句。]

我们考虑某个固定的拓扑空间 $X$ 作底空间；我们想研究 $X$ 上所有有限维（实或复）向量丛构成的集合 $VB(X)$。在这个集合上可以定义下列运算：

\begin{enumerate}[label=\arabic*.]
\item \emph{直和}（direct sum）：若 $\xi,\eta$ 是 $X$ 上的两个向量丛，我们可以定义它们的
% [原文 p.14]
直和 $\xi\oplus\eta$，它是 $X$ 上的一个新向量丛，其全空间为
\[ E(\xi\oplus\eta)\subset E(\xi)\times E(\eta) \]
其中
\[ E(\xi\oplus\eta)=\{(v,w) : \pi_\xi(v)=\pi_\eta(w)\}. \]
这个新向量丛的维数等于 $\xi,\eta$ 的维数之和，并且每条纤维是对应纤维的直和（回忆纤维是向量空间，因此可以取它们的直和）
\[ F_x(\xi\oplus\eta)=F_x(\xi)\oplus F_x(\eta). \]

\item \emph{拉回}（pull-back）：若 $f: X\to Y$ 是连续映射，$\xi\to Y$ 是 $Y$ 上的向量丛，则可以定义向量丛 $f^*(\xi)\to X$，称为 $\xi$ 沿 $f$ 在 $X$ 上的\emph{拉回}，它是 $X$ 上一个新的向量丛，其全空间为
\[ E(f^*(\xi))\subset X\times E(\xi) \]
其中
\[ E(f^*(\xi))=\{(x,e) : \pi_\xi(e)=f(x)\}. \]
它的纤维定义如下：
\[ F_x(f^*(\xi))=F_{f(x)}(\xi). \]

\item \emph{张量积}（tensor product）：设 $\pi_1: E_1\to X$ 与 $\pi_2: E_2\to X$ 是同一底空间 $X$ 上的两个向量丛。则我们可以构造它们的张量积 $E_1\otimes E_2$，它是 $X$ 上的一个向量丛。作为集合，全空间 $E_1\otimes E_2$ 是诸（纤维）向量空间张量积 $\pi_1^{-1}(x)\otimes\pi_2^{-1}(x)$ 的不交并。$X$ 上向量丛的结构可以用粘合函数给出：取 $X$ 的开覆盖 $\{U_a\}$，使得 $E_i$（$i=1,2$）在每个 $U_a$ 上平凡，于是对每个 $E_i$ 得到粘合函数 $g^i_{ba}: U_a\cap U_b\to GL_{n_i}(\F)$。则张量丛 $E_1\otimes E_2$ 的粘合函数由张量积函数 $g^1_{ba}\otimes g^2_{ba}$ 给出，后者对每个 $x\in U_a\cap U_b$ 指派两个矩阵 $g^1_{ba}(x)$ 与 $g^2_{ba}(x)$ 的张量积。

% [Tutor Note: p.14 张量丛的纤维维数是 $n_1n_2$：$g^1_{ba}\otimes g^2_{ba}$ 是 $GL_{n_1}(\F)\times GL_{n_2}(\F)\to GL_{n_1n_2}(\F)$ 的张量积表示（即矩阵的 Kronecker 积），故 $E_1\otimes E_2$ 的秩为 $\rk E_1\cdot\rk E_2$。原文只写粘合函数，未提维数。]

\item \emph{商丛}（quotient bundles）：给定向量丛 $\pi: E\to X$ 及 $E$ 的向量子丛 $E_0\subset E$（它本身是 $X$ 上的另一个向量丛），我们可以
% [原文 p.15]
构造商丛 $E/E_0\to X$，其全空间带商拓扑，而任意点 $x\in X$ 上的纤维是对应纤维的商 $F_x(E)/F_x(E_0)$；显然商丛 $E/E_0$ 的维数等于 $E$ 与 $E_0$ 的维数之差，即
\[ \dim(E/E_0)=\dim E-\dim E_0. \]
\end{enumerate}

\para{习题 1.} \emph{定义向量丛的外积。}

\para{习题 2.} \emph{证明拉回同时保持直和与张量积，即 $f^*(E_1\oplus E_2)=f^*(E_1)\oplus f^*(E_2)$ 且 $f^*(E_1\otimes E_2)=f^*(E_1)\otimes f^*(E_2)$。}

显然，平凡向量丛的直和仍是平凡向量丛；一个非平凡向量丛若与某个平凡向量丛作直和后变成平凡丛，则称为\emph{稳定平凡}（stably trivial）；不那么显然的是，非平凡向量丛的直和可以是平凡的；在一定假设下，后一种情况总是会发生：

\begin{prop}{1}
对任意向量丛 $\pi: E\to B$，其中 $B$ 紧且 Hausdorff，存在另一个向量丛 $E'\to B$，使得 $E\oplus E'$ 是平凡的。
\end{prop}

为了证明这个命题，我们需要一个定义与一条引理：

\begin{defn}{7}
向量丛 $\pi: E\to B$ 上的\emph{内积}（inner product）是一个映射 $\langle\,,\rangle: E\oplus E\to\R$，它在每条纤维上的限制是一个内积（即一个正定对称双线性型）。
\end{defn}

［注：为简短起见，这里只处理实情形；复情形也同样可以直接得到。］

% [Tutor Note: p.15 方括号注中原文 "compelx" 是 "complex" 的笔误。另定义 7 把内积写成 $E\oplus E\to\R$，此处 $E\oplus E$ 指纤维逐点取直和所得的丛（见定义 5），而非拓扑积 $E\times E$。]

当底空间紧且 Hausdorff（更一般地，仿紧——即这样的拓扑空间：它 Hausdorff，且每个开覆盖都有从属于它的单位分解，而单位分解可用 Urysohn 引理构造）时，内积总存在。内积可以这样得到：把 $\R^n$ 中的标准欧氏内积经局部平凡化拉回，
% [原文 p.16]
然后用单位分解把它延拓到整个 $B$ 上。

\begin{lem}{1}
若 $\pi: E\to B$ 是仿紧空间 $B$ 上的向量丛，$E_0\subset E$ 是向量子丛，则存在另一个向量子丛 $E_0^\perp\subset E$，使得 $E_0\oplus E_0^\perp\simeq E$。
\end{lem}

\begin{proof}[引理 1 的证明]
设 $E$ 的维数为 $n$，$E_0$ 的维数为 $m<n$。我们在 $E$ 上取一个内积，并定义 $E_0^\perp$ 为 $E$ 的这样的子空间：其纤维由与 $E_0$ 的向量正交的向量组成。我们断言自然投影 $E_0^\perp\to B$ 定义 $B$ 上的一个向量丛（从而使 $E_0^\perp$ 成为子丛，而不只是 $E$ 的子空间）。若果然如此，我们就得到所需结论，即 $E\simeq E_0\oplus E_0^\perp$，同构由 $(v,w)\mapsto v+w$ 给出（因为向量丛的同构限制到每条纤维上是向量空间的线性同构）。

我们必须验证 $E_0^\perp$ 上的局部平凡性条件。由于 $E_0$ 的维数为 $m$，它在每点 $b_0\in B$ 附近有 $m$ 个线性无关的局部截面 $b\mapsto(b,s_i(b))$。我们把 $E_0$ 的这 $m$ 个线性无关局部截面扩充为 $E$ 的 $n$ 个线性无关局部截面：先在纤维 $\pi^{-1}(b_0)$ 中选取 $s_{m+1},\dots,s_n$，然后对附近所有纤维取同样的向量，因为若 $s_1,\dots,s_m,\dots,s_n$ 在 $b_0$ 处线性无关，则由行列式的连续性，它们在附近的 $b$ 处仍线性无关。对 $s_1,\dots,s_m,\dots,s_n$ 施行 Gram--Schmidt 正交化过程（使用每条纤维中的内积），得到新的截面 $s_i'$。正交化过程的显式公式表明新的局部截面连续，且其中前 $m$ 个在每条纤维上给出 $E_0$ 的一组基。新的截面 $s_i'$ 使我们能定义局部平凡化 $h: \pi^{-1}(U)\to U\times\R^n$，使得 $h(b,s_i'(b))$ 等于 $\R^n$ 的第 $i$ 个标准基向量。这个 $h$ 把 $E_0$ 送到 $U\times\R^m$，把 $E_0^\perp$ 送到 $U\times\R^{n-m}$，因此 $h|_{E_0^\perp}$ 是 $E_0^\perp$ 的局部平凡化。
\end{proof}

% [Tutor Note: p.16 引理 1 证明的关键补白：把 $m$ 个截面延拓为 $n$ 个时，「对附近纤维取同样的向量」只在 $b_0$ 的一个小邻域内有效，正是行列式的连续性保证该邻域存在；随后 Gram--Schmidt 的过渡矩阵由 $s_i$ 连续可逆地给出，故 $s_i'$ 连续。原文「$h$ 把 $E_0^\perp$ 送到 $U\times\R^{n-m}$」可由纤维维数逐点核对。另原文把 "isomorphisms" 拼作 "isomorhisms"（括号内 "since isomorhisms of vector bundles restrict..."）。]

\begin{proof}[命题 1 的证明]
为了说明这一构造的动机，先假定结论成立，于是 $E$ 是平凡丛 $B\times\R^N$（$N$ 足够大）的子丛，从而有一个包含 $E\into B\times\R^N$；把这一包含与到第二个因子的投影复合，得到映射 $E\to\R^N$，它在每条纤维上是线性单射；证明的想法是把逻辑倒过来，即先构造一个在每条纤维上为线性单射的映射 $E\to\R^N$，然后证明这给出 $E$ 在 $B\times\R^N$ 中作为一个直和项的嵌入。

% [原文 p.17]
$B$ 的每一点 $x$ 都有一个邻域 $U_x$，使 $E$ 在其上平凡。由 Urysohn 引理，存在映射 $\phi_x: B\to[0,1]$，它在 $U_x$ 之外为零、在 $x$ 处非零。令 $x$ 变动，集合 $\phi_x^{-1}(0,1]$ 构成 $B$ 的开覆盖；又因 $B$ 紧，该覆盖有有限子覆盖。我们把相应的 $U_x'$ 与 $\phi_x'$ 分别重新标号为 $U_j$ 与 $\phi_j$。定义 $g_j: E\to\R^n$ 为
\[ g_j(v)=\phi_j(\pi(v))[p_jh_j(v)], \]
其中 $p_jh_j$ 是局部平凡化 $h_j: \pi^{-1}(U_j)\to U_j\times\R^n$ 与到 $\R^n$ 的投影 $p_j$ 的复合。则 $g_j$ 在 $\phi_j^{-1}(0,1]$ 上的每条纤维上是线性单射；于是，若把诸 $g_j$ 取作映射 $g: E\to\R^N$ 的坐标（其中 $\R^N$ 是若干 $\R^n$ 的乘积），则 $g$ 在每条纤维上是线性单射。

映射 $g$ 是映射 $f: E\to B\times\R^N$ 的第二个坐标，其第一个坐标是 $\pi$。$f$ 的像是乘积 $B\times\R^N$ 的子丛，因为把 $\R^N$ 投影到第 $i$ 个 $\R^n$ 因子所得恰是 $\phi_j^{-1}(0,1]$ 上一个局部平凡化的第二个坐标。于是 $E$ 同构于 $B\times\R^N$ 的一个子丛，故由引理 1，存在补子丛 $E'$ 使 $E\oplus E'$ 同构于 $B\times\R^N$。
\end{proof}

% [Tutor Note: p.17 构造 $g_j$ 时乘上 $\phi_j(\pi(v))$ 是为了让 $g_j$ 在 $U_j$ 外自动为零（从而能作为整体映射的坐标）；而「$g_j$ 是纤维单射」只在 $\phi_j^{-1}(0,1]$ 上成立，因为那里 $\phi_j\neq0$ 只是标量因子，不改变线性单射性。核对时可逐纤维看。]

\emph{下文我们总假定拓扑空间 $X$ 紧且 Hausdorff，且所讨论的向量丛是复的。}

设 $\Vect_k(X)$ 表示 $X$ 上 $k$ 维向量丛的同构类集合，$\Vect(X)$ 表示 $X$ 上有限维向量丛的同构类集合，$VB(X)$ 表示 $X$ 上所有有限维向量丛的集合。则上面定义的直和 $\oplus$ 在 $\Vect(X)$ 上定义一个加法：
\[ [\xi]+[\eta]=[\xi\oplus\eta]. \]
这个运算是结合的、交换的，存在中性元（零元），但无法定义对称元（丛的「相反」同构类），因此 $(\Vect(X),+)$ 是一个（加法阿贝尔）\emph{半群}（semi-group）。

有一个把半群凑成群的技巧，可追溯到 Grothendieck：这是一种对称化过程，正如我们由（加法）半群 $\N$ 定义 $\Z$ 一样。更具体地说，对
% [原文 p.18]
任意半群 $M$，我们为其联系一个新群，记作 $Gr(M)$，它定义为 $M\times M$ 按等价关系
\[ (a,b)\sim(c,d)\iff\exists e\in M : a+d+e=b+c+e \]
作成的商。

% [Tutor Note: p.18 等价关系中多出的 $e$ 是必需的：半群没有消去律，不能直接令 $a+d=b+c$；添上公共的 $e$ 才使 $\sim$ 传递——若 $(a,b)\sim(c,d)$ 且 $(c,d)\sim(f,g)$，则取 $e+e'$ 作公共项即得 $(a,b)\sim(f,g)$。这是 Grothendieck 群构造的标准细节。]

还有一种做同样事情的替代方法：若 $M$ 是任意半群，则考虑半群的对角同态
\[ \Delta: M\to M\times M. \]
则 $Gr(M)$（半群 $M$ 的 Grothendieck 群或万有群）是 $\Delta(M)$ 在 $M\times M$ 中的陪集集合，它是商半群；然而笛卡尔积 $M\times M$ 中因子的互换在 $Gr(M)$ 中诱导出对称元，于是它被提升为一个群。

回到我们的 K 理论，于是有了下面这个归于 Alexander Grothendieck 的关键定义：

\begin{defn}{8}
我们定义阿贝尔群 $K(X)$ 为半群 $\Vect(X)$ 的万有群（或称「Grothendieck 群」，记作 $Gr$），即
\[ K(X)=Gr[\Vect(X)]. \]
（事实上这就是拓扑 K 理论的第 0 个 K 群，我们后面会看到。）［通常把复情形与实情形分别记作 $K(X)$ 与 $KO(X)$］。
\end{defn}

\begin{defn}{9}
同一底空间 $X$ 上的两个向量丛 $E_1$ 与 $E_2$ 称为\emph{稳定同构的}（stably isomorphic），记作 $E_1\simeq_S E_2$，如果存在某个 $n\in\N$ 使得 $E_1\oplus\C^n_X\simeq E_2\oplus\C^n_X$。
\end{defn}

\begin{defn}{10}
同一底空间 $X$ 上的两个向量丛 $E_1$ 与 $E_2$ 称为\emph{相似的}（similar），记作 $E_1\sim E_2$，如果存在某些 $n,m\in\N$ 使得 $E_1\oplus\C^n_X\simeq E_2\oplus\C^m_X$。
\end{defn}

% [Tutor Note: p.18 定义 9 与定义 10 的差别只在平凡丛的维数：稳定同构要求两边加同一个 $\C^n_X$，相似则允许两边加不同维数的平凡丛。于是 $\simeq_S$ 更细（稳定同构 $\Rightarrow$ 相似），而 $\sim$ 的类才构成命题 2 中的群 $\wt K(X)$。另定义 8 之后方括号内原文 "Its customarily to denote KX and KOX and the complex and the real case respectively" 中 "Its" 应为 "It is"、且并列语序混乱；译文按原意理顺为「复情形与实情形分别记作 $K(X)$ 与 $KO(X)$」。]

证明 $\simeq_S$ 与 $\sim$ 都是 $VB(X)$ 中的等价关系，确实十分直接。在两种等价类上直和都良定义，而且它是结合、交换的。存在中性（零）元，
% [原文 p.19]
即 $\C^0_X$ 的类。

\begin{defn}{8'}
阿贝尔群 $K(X)$ 的一个等价定义是：$K(X)$ 是 $VB(X)$ 按\emph{稳定同构}这一等价关系作成的商空间，即
\[ K(X)=VB(X)/\simeq_S. \]
\end{defn}

% [Tutor Note: p.19 定义 8' 与定义 8 的等价性依赖命题 1（紧 Hausdorff 底空间上每个丛都有补丛）：定义 8 用形式差 $[E]-[E']$ 造群，定义 8' 直接用稳定同构类，二者在紧 Hausdorff 且纤维维数一致时给出同一个群。]

于是「稳定性」正是由半群得到阿贝尔群的 Grothendieck 技巧的编码。

\begin{prop}{2}
若底空间 $X$ 连通、紧且 Hausdorff，则相似关系下的等价类集合 $VB(X)/\sim$ 在直和 $\oplus$ 下构成另一个阿贝尔群，记作 $\wt K(X)$。
\end{prop}

\begin{defn}{11}
群 $\wt K(X)$ 称为 $X$ 的\emph{约化 K 理论}（$X$ 的约化第 0 个 K 群）。
\end{defn}

\begin{proof}[命题 2 的证明]
只需证明逆元存在；这通过说明对每个向量丛 $\pi: E\to X$ 都存在向量丛 $E'\to X$ 使得 $E\oplus E'\simeq\C^n_X$（对某个 $n$）来完成。而这恰恰就是上面的命题 1。（我们假定了 $X$ 连通，从而所有纤维维数相同；不难证明即使放松连通性假设，命题 2 仍成立。）
\end{proof}

% [Tutor Note: p.19 命题 2 的证明只用到了命题 1，而命题 1 要求底空间紧 Hausdorff；括注里说「放松连通性也不难」，但紧 Hausdorff 这一条不能去掉——去掉后须另辟途径（例如改用 $C(X)$ 的投射模）才能定义 K 群。]

若把向量丛之间的张量积作为额外的运算，则 $\Vect(X)$ 成为一个半环，而 $K(X)$ 是一个环。

对于稳定同构的 K 类上的直和运算，只有零元 $[\C^0_X]$ 能有逆元，因为 $E\oplus E'\simeq_S\C^0_X$ 蕴涵 $E\oplus E'\oplus\C^n_X\simeq_S\C^n_X$（对某个 $n$），而这只有当 $E$ 与 $E'$ 都是零维时才可能发生。然而，尽管逆元不存在，我们确实有消去性质：在紧空间 $X$ 上 $E_1\oplus E_2\simeq_S E_1\oplus E_3\Rightarrow E_2\simeq_S E_3$，因为可以给 $E_1\oplus E_2\simeq_S E_1\oplus E_3$ 两边加上一个丛 $E_1'$，使得 $E_1\oplus E_1'\simeq\C^n_X$（对某个 $n$）。

于是对紧的 $X$，K 类可以表示成向量丛的「形式差」$E-E'$，其中 $E$ 与 $E'$ 是 $X$ 上真正的实向量丛，
% [原文 p.20]
它们带有等价关系
\[ E_1-E_1'=E_2-E_2' \]
当且仅当
\[ E_1\oplus E_2'\simeq_S E_2\oplus E_1'. \]
为了使这一条件良定义，我们需要消去性质，从而 $X$ 必须紧。按显然的加法规则
\[ (E_1-E_1')+(E_2-E_2')=(E_1\oplus E_2)-(E_1'\oplus E_2') \]
$K(X)$ 成为一个群。零元是任意 $E$ 对应的 $E-E'$ 的等价类，而 $E-E'$ 的逆元是 $E'-E$。此外注意，$K(X)$ 的每个元素也可以表示成差
\[ E-\C^n_X, \]
因为若从 $E-E'$ 出发，可以给 $E$ 与 $E'$ 都加上一个丛 $E''$，使得 $E'\oplus E''\simeq\C^n_X$（对某个 $n$）。

% [Tutor Note: p.19 末句原文 "honest real vector bundles"（真正的实向量丛）与前文「本章向量丛取复的」这一约定不一致，疑为笔误（应泛指复向量丛，或作者沿用了实情形的说法）；等价关系的定义本身与 $\R$ 或 $\C$ 无关，故不影响结论。]

存在自然的同态 $K(X)\to\wt KX$，它把 $E-\C^n_X$ 送到 $E$ 对应的 $\sim$-类；这是良定义的，因为若在 $K(X)$ 中 $E-\C^n_X=E'-\C^m_X$，则 $E\oplus\C^m_X\simeq_S E'\oplus\C^n_X$，从而 $E\sim E'$。这个映射显然满，其核由满足 $E\sim\C^0_X$ 的元素 $E-\C^n_X$ 组成，于是 $E\simeq_S\C^m_X$（对某个 $m$），故核由形如 $\C^m_X-\C^n_X$ 的元素组成。$K(X)$ 的这个子群同构于 $\Z$。事实上，把向量丛限制到基点 $x_0\in X$ 定义了同态 $K(X)\to K(x_0)\simeq\Z$，它限制在子群 $\{\C^m_X-\C^n_X\}$ 上给出同构。于是有分裂
\[ K(X)\simeq\wt KX\oplus\Z, \]
它依赖于 $x_0$ 的选取。

% [Tutor Note: p.20 核同构于 $\Z$ 的核对方式：限制到基点给出 $K(X)\to K(x_0)\simeq\Z$，而 $\C^m_X-\C^n_X$ 被送到 $m-n$，可见该子群同构于 $\Z$；分裂 $K(X)\simeq\wt KX\oplus\Z$ 依赖基点 $x_0$ 的选取，原文特别指出这一点。]

设 $f: X\to Y$ 是连续映射，$E\to Y$ 是 $Y$ 上的向量丛。则如前所见，可以构造 $E$ 沿 $f$ 在 $X$ 上的拉回丛 $f^*(E)$。于是任何这样的映射 $f: X\to Y$ 诱导映射 $f^*: \Vect(Y)\to\Vect(X)$，其中 $\xi\mapsto f^*\xi$，它保持直和，即是半群同态，从而给出阿贝尔群之间的群同态
% [原文 p.21]
$f^*: K(Y)\to K(X)$。这常被称为 K 理论中的\emph{错误方向函子性}（the wrong way functoriality in K-Theory）。

\begin{thm}{1}
设 $f_0,f_1: X\to Y$ 是同伦映射，$\xi=(E,\pi,Y)$ 是 $Y$ 上的向量丛。则 $f_0^*\xi$ 与 $f_1^*\xi$ 同构。
\end{thm}

回忆 $f_0,f_1: X\to Y$ 同伦是指存在连续映射 $F: X\times I\to Y$（其中 $I=[0,1]$），使得 $f_0(x)=F(x,0)$ 且 $f_1(x)=F(x,1)$。

\para{定理 1 的证明}
设 $F: X\times I\to Y$ 是从 $f_0$ 到 $f_1$ 的同伦。$F^*(E)$ 在 $X\times\{0\}$ 与 $X\times\{1\}$ 上的限制分别是 $f_0^*\xi$ 与 $f_1^*\xi$。因此只需证明下面这一点：

% [Tutor Note: p.21 原文的「Proof of Theorem 1:」与随后引理 2 的证明共用末尾一个 □（原文只在引理 2 证明末给出一个证毕符号），故此处定理 1 的证明只写标题、不封口：其证明由引理 2 及其证明补全。译文照此。]

\begin{lem}{2}
向量丛 $E\to X\times I$ 在 $X\times\{0\}$ 与 $X\times\{1\}$ 上的限制同构，只要 $X$ 紧且 Hausdorff（更一般地，只要 $X$ 仿紧）。
\end{lem}

\begin{proof}[引理 2 的证明]
我们将用到两个预备事实：

\begin{enumerate}[label=\roman*.]
\item 向量丛 $\pi: E\to X\times[a,b]$ 是平凡的，如果它在 $X\times[a,c]$ 与 $X\times[c,b]$ 上的限制对某个 $c\in(a,b)$ 都平凡。为看出这一点，设这两个限制为 $E_1=\pi^{-1}(X\times[a,c])$ 与 $E_2=\pi^{-1}(X\times[c,b])$，并设 $h_1: E_1\to X\times[a,c]\times\C^n$ 与 $h_2: E_2\to X\times[c,b]\times\C^n$ 是同构。这些同构在 $\pi^{-1}(X\times\{c\})$ 上可能不一致，但可以通过把 $h_2$ 复合上一个同构来使之一致：该同构 $X\times[c,b]\times\C^n\to X\times[c,b]\times\C^n$ 在每个切片 $X\times\{x\}\times\C^n$ 上由 $h_1h_2^{-1}: X\times\{c\}\times\C^n\to X\times\{c\}\times\C^n$ 给出。一旦 $h_1$ 与 $h_2$ 在 $E_1\cap E_2$ 上一致，它们就定义了 $E$ 的一个平凡化。

% [Tutor Note: p.21 事实 (i) 中原文 "on each slice $X\times\{x\}\times\C^n$" 疑误：切片应由 $t=c$ 参数化，应为 $X\times\{c\}\times\C^n$（$x$ 是 $X$ 中的点，不能同时充当切片参数）。同句的 $h_1h_2^{-1}$ 正定义在 $X\times\{c\}\times\C^n$ 上，可据此核对。]

\item 对向量丛 $\pi: E\to X\times I$，存在 $X$ 的开覆盖 $\{U_a\}$，使得每个限制 $\pi^{-1}(U_a\times I)\to U_a\times I$ 平凡。这是因为对每个 $x\in X$ 可以找到 $X$ 中的开邻域 $U_{x,1},\dots,U_{x,k}$ 以及 $[0,1]$ 的一个分划 $0=t_0<t_1<\dots<t_k=1$，使得丛在 $U_{x,i}\times[t_{i-1},t_i]$ 上平凡（这里用到 $[0,1]$ 的紧性）。于是由事实 (i)，丛在 $U_a\times I$ 上平凡，其中 $U_a=U_{x,1}\cap\dots\cap U_{x,k}$。
\end{enumerate}

现在我们来证明引理 2：由事实 (ii) 我们可以选取
% [原文 p.22]
$X$ 的开覆盖 $\{U_a\}$，使得 $E$ 在每个 $U_a\times I$ 上平凡。先处理 $X$ 紧 Hausdorff 这一较简单的情形：此时有限个 $U_a$ 覆盖 $X$。把它们重新标号为 $U_i$，$i=1,2,\dots,m$。如命题 1 中所说，存在相应的单位分解，由支撑包含在 $U_i$ 中的函数 $\phi_i$ 组成。对 $i\ge1$，令 $\psi_i=\phi_1+\dots+\phi_i$，特别地 $\psi_0=0$、$\psi_m=1$。设 $X_i$ 为 $\psi_i$ 的图像，即 $X\times I$ 中由形如 $(x,\psi_i(x))$ 的点组成的子空间；并设 $\pi_i: E_i\to X_i$ 为 $E$ 在 $X_i$ 上的限制。由于 $E$ 在 $U_i\times I$ 上平凡，自然投影同胚 $X_i\to X_{i-1}$ 可提升为同胚 $h_i: E_i\to E_{i-1}$，它在 $\pi_i^{-1}(U_i)$ 之外是恒等映射，并把 $E_i$ 的每条纤维同构地映到 $E_{i-1}$ 的对应纤维上。于是复合 $h=h_1h_2\dots h_m$ 是从 $E$ 在 $X\times\{1\}$ 上的限制到其在 $X\times\{0\}$ 上的限制的同构。

% [Tutor Note: p.22 复合顺序 $h=h_1h_2\dots h_m$ 要按箭头走向读：$h_i: E_i\to E_{i-1}$，故从 $E_m$（在 $X\times\{1\}$ 上）出发依次经 $h_m,\dots,h_1$ 落到 $E_0$（在 $X\times\{0\}$ 上）；原文按指标升序写成 $h_1h_2\dots h_m$，按通常的函数复合（右作用）顺序理解时须注意这一点。]

在 $X$ 只是仿紧的一般情形，存在 $X$ 的可数覆盖 $\{V_i\}_{i\ge1}$ 与单位分解 $\{\phi_i\}$，其中 $\phi_i$ 支撑在 $V_i$ 中，且每个 $V_i$ 是一些开集的不交并，每个开集都包含在某个 $U_a$ 中。这意味着 $E$ 在每个 $V_i\times I$ 上平凡。如前令 $\psi_i=\phi_1+\dots+\phi_i$，并令 $\pi_i: E_i\to X_i$ 为 $E$ 在 $\psi_i$ 的 grapf（原文如此，应为 graph，即图像）上的限制；我们利用 $E$ 在 $V_i\times I$ 上平凡这一事实构造 $h_i: E_i\to E_{i-1}$。于是无限复合 $h=h_1h_2\dots$ 是从 $E$ 在 $X\times\{1\}$ 上的限制到其在 $X\times\{0\}$ 上的限制的良定义同构，因为在每点 $x\in X$ 附近只有有限多个 $\phi_i$ 非零，从而存在 $x$ 的一个邻域，其中除有限多个 $h_i$ 外其余都是恒等映射。
\end{proof}

\begin{cor}{1}
若 $X$ 可缩（即同伦于一点），则 $\Vect(X)=\N$ 且 $K(X)=\Z$。
\end{cor}

\begin{proof}
$X$ 可缩意味着它同伦于一点，因此一点上的向量丛不过是一个向量空间；向量空间的同构类由其维数刻画，故 $\Vect(*)=\N$，于是万有（即 Grothendieck）群 $Gr(\N)=\Z$。
\end{proof}

\begin{cor}{2}
可缩空间上的向量丛是平凡的。
\end{cor}

\begin{proof}
直接。
\end{proof}

% [Tutor Note: p.22 推论 1 证明中 $\Vect(*)=\N$ 的「$*$」表示单点空间；由引理 2 的同伦不变性，可缩空间 $X$ 上的丛与单点上的丛对应，故 $K(X)\simeq K(*)\simeq\Z$。原文把这一步压缩成一句。]

% [原文 p.23]
［定理 1 对实向量丛同样成立。它对任意纤维丛也成立］。

我们定义了拓扑空间上的向量丛，并定义了它们之间的各种等价关系（同构、稳定同构、相似）。一个重要的问题是某个固定拓扑空间（比如 $X$）上向量丛的\emph{分类}：相似意义下的分类可以借助约化 K 理论实现，稳定同构意义下的分类可以借助 K 理论完成，但同构意义下的分类问题在很大程度上仍是公开问题。处理这个问题的一个非常有用的工具是一些称为「\emph{示性类}」（characteristic classes）的上同调对象，我们后面会遇到它们。它们给出了向量丛同构分类问题的部分答案；事实上在低维情形它们非常好用。眼下我们将用代数拓扑的一个标准概念——映射的同伦类的想法，连同分类空间与万有丛的概念——来重新表述分类问题。于是我们将构造一个 $k$ 维向量丛（称为万有丛）$E_k\to G_k$，它具有这样的性质：紧 Hausdorff 空间上所有 $k$ 维向量丛都可以作为这个单一丛的拉回得到（特别地，这可推广到仿紧空间；当 $k=1$ 时，在实情形下给出无限射影空间上的线丛）。

回忆 $G_k(\C^n)$ 表示 $\C^n$ 的 $k$ 维向量子空间的 Grassmann 流形（$k\le n$ 均为非负整数，即过原点的 $k$ 维平面的集合）。可以借助 Stiefel 流形 $V_k(\C^n)$ 给 Grassmann 流形赋予拓扑，后者是 $\C^n$ 中正交 $k$-标架的空间，换言之是 $\C^n$ 中正交向量构成的 $k$-元组。它是 $k$ 个单位球面 $S^{n-1}$ 的乘积的子空间，即正交 $k$-元组构成的子空间。它是闭子空间，因为两个向量的正交性可以用一个代数方程表达。因此 Stiefel 流形紧（因为球面的乘积紧）。存在自然满射 $V_k(\C^n)\to G_k(\C^n)$，把一个 $k$-标架送到它张成的子空间；于是可以赋予 $G_k(\C^n)$ 关于这个满射的商拓扑，从而 Grassmann 流形也紧。

% [Tutor Note: p.23 原文 "One can gine a topology" 中 "gine" 是 "give" 的笔误；同段 "the k-tuples of of orthonormal vectors" 多写了一个 "of"。另「$V_k(\C^n)$ 是 $k$ 个 $S^{n-1}$ 的乘积的子空间」精确说法是：它由正交方程组（代数方程）在 $(S^{n-1})^k$ 中切出，故闭。]

包含关系 $\C^n\subset\C^{n+1}\subset\dots$ 给出包含关系 $G_k(\C^n)\subset G_k(\C^{n+1})\subset\dots$
% [原文 p.24]
并且我们令
\[ G_k(\C^\infty)=\bigcup_n G_k(\C^n). \]
我们给 $G_k(\C^\infty)$ 赋以弱（或正向极限）拓扑：一个集合是开集当且仅当它与每个 $G_k(\C^n)$ 的交都是开集。

在 $G_k(\C^n)$ 上有典范的 $n$ 维向量丛：定义
\[ E_k(\C^n)=\{(l,v)\in G_k(\C^n)\times\C^n : v\in l\}. \]

% [Tutor Note: p.24 原文说 $G_k(\C^n)$ 上有「典范的 $n$ 维向量丛」，但 $E_k(\C^n)$ 的纤维是 $v\in l$（$l$ 为 $k$ 维子空间），维数应为 $k$；由下文引理 3 及 $E_k\to G_k$ 是 $k$ 维丛的分类空间可知这里的 "n-dim" 应为 "k-dim"。译文照原样作「$n$ 维」。同页原文 "The condtition on $v_1,\dots,v_k$" 中 "condtition" 是 "condition" 的笔误。]

包含关系 $\C^n\subset\C^{n+1}\subset\dots$ 给出包含关系 $E_k(\C^n)\subset E_k(\C^{n+1})\subset\dots$，并令
\[ E_k(\C^\infty)=\bigcup_n E_k(\C^n), \]
同样带正向极限拓扑。

\begin{lem}{3}
投影 $p: E_k(\C^n)\to G_k(\C^n)$（$p(l,v)=l$）对有限的 $n$ 与无限的 $n$ 都定义一个向量丛。
\end{lem}

\begin{proof}
先设 $n$ 有限。对 $l\in G_k(\C^n)$，设 $\pi_l: \C^n\to l$ 是正交投影，并令 $U_l=\{l'\in G_k(\C^n) : \dim\pi_l(l')=k\}$。特别地 $l\in U_l$。我们将证明 $U_l$ 在 $G_k(\C^n)$ 中开，且映射 $h: p^{-1}(U_l)\to U_l\times l\simeq U_l\times\C^k$（定义为 $h(l',v)=(l',\pi_l(v))$）是 $E_k(\C^n)$ 的局部平凡化。

$U_l$ 开等价于它在 $V_k(\C^n)$ 中的原像开。该原像由正交标架 $v_1,\dots,v_k$ 组成，使得 $\pi_l(v_1),\dots,\pi_l(v_k)$ 线性无关。设 $A$ 是 $\pi_l$ 关于定义域 $\C^n$ 中标准基与值域 $l$ 中某个固定基的矩阵。则关于 $v_1,\dots,v_k$ 的条件是：以 $Av_1,\dots,Av_k$ 为列的 $k\times k$ 矩阵行列式非零。由于该行列式的值显然是 $v_1,\dots,v_k$ 的连续函数，可知使行列式非零的标架 $v_1,\dots,v_k$ 构成 $V_k(\C^n)$ 中的开集。

显然 $h$ 是双射，且在每条纤维上是线性同构。我们需要验证 $h$ 与 $h^{-1}$ 连续。对 $l'\in U_l$，存在唯一的可逆线性映射 $L_{l'}: \C^n\to\C^n$，它在 $l'$ 上限制为 $\pi_l$、在 $l^\perp=\Ker(\pi_l)$ 上是恒等映射。我们断言，把 $L_{l'}$ 看作 $n\times n$ 矩阵时它
% [原文 p.25]
连续地依赖于 $l'$，即可以写成
\[ L_{l'}=AB^{-1}, \]
其中 $B$ 把标准基送到 $v_1,\dots,v_k,v_{k+1},\dots,v_n$，这里 $v_1,\dots,v_k$ 是 $l'$ 的一组正交基，$v_{k+1},\dots,v_n$ 是 $l^\perp$ 的一组固定基；而 $A$ 把标准基送到 $\pi_l(v_1),\dots,\pi_l(v_k),v_{k+1},\dots,v_n$。

$A$ 与 $B$ 都连续依赖于 $v_1,\dots,v_k$。由于矩阵乘法与矩阵求逆都是连续运算，可知乘积 $L_{l'}=AB^{-1}$ 连续依赖于 $v_1,\dots,v_k$。又因 $L_{l'}$ 只依赖于 $l'$ 而不依赖于 $l'$ 的基 $v_1,\dots,v_k$，故 $L_{l'}$ 连续依赖于 $l'$（因为 $G_k(\C^n)$ 带 $V_k(\C^n)$ 的商拓扑）。由于 $h(l',v)=(l',\pi_l(v))=(l',L_{l'}(v))$，可见 $h$ 连续。类似地 $h^{-1}(l',w)=(l',L_{l'}^{-1}(w))$，而 $L_{l'}^{-1}$ 连续依赖于 $l'$（矩阵求逆连续），故 $h^{-1}$ 连续。这就完成了有限 $n$ 情形的证明。

对无限情形，取 $U_l$ 为递增的 $n$ 所对应的诸 $U_l$ 的并。上面为有限情形构造的局部平凡化 $h$ 便拼合起来给出这个 $U_l$ 上的局部平凡化；而连续性显然，因为我们用的是弱拓扑。
\end{proof}

% [Tutor Note: p.24--25 引理 3 证明中 $L_{l'}$ 的定义是关键：它在 $l'$ 上等于 $\pi_l$、在 $l^\perp$ 上等于恒等；$l'\in U_l$ 意味着 $\dim\pi_l(l')=k$，故 $\pi_l$ 在 $l'$ 上单射，$L_{l'}$ 可逆。局部平凡化的性质正来自 $h(l',v)=(l',L_{l'}(v))$ 逐纤维为线性同构。]

我们现在关心 $n=\infty$ 的情形，为简化记号，写 $G_k$ 表示 $G_k(\C^\infty)$，类似地写 $E_k$ 表示 $E_k(\C^\infty)$。如前所述，用 $[X,Y]$ 记映射 $f: X\to Y$ 的同伦类的集合。

\begin{thm}{2}
对紧 Hausdorff 的 $X$（更一般地，若 $X$ 仿紧），映射 $[X,G_k]\to\Vect_k(X)$，$[f]\mapsto f^*(E_k)$，是双射。
\end{thm}

\begin{proof}
关键的观察如下：对 $k$ 维向量丛 $\pi: E\to X$，同构 $E\simeq f^*(E_k)$ 等价于一个在每条纤维上为线性单射的映射 $g: E\to\C^\infty$。为看出这一点，先设有一个映射 $f: X\to G_k$ 以及同构 $E\simeq f^*(E_k)$。于是有
% [原文 p.26]
交换图
\[
\begin{tikzcd}
E \arrow[r,"\simeq"] \arrow[d,"p"'] & f^*(E_k) \arrow[r,"\tilde f"] \arrow[d] & E_k \arrow[r,"\pi"] \arrow[d] & \C^\infty \\
X \arrow[r,"id"'] & X \arrow[r,"f"'] & G_k &
\end{tikzcd}
\]
其中 $\pi(l,v)=v$。上排的复合是映射 $g: E\to\C^\infty$，它在每条纤维上是线性单射，因为 $\tilde f$ 与 $\pi$ 都有这一性质。反过来，给定在每条纤维上是线性单射的映射 $g: E\to\C^\infty$，令 $f(x)$ 为 $k$ 维平面 $g(p^{-1}(x))$，便定义了 $f: X\to G_k$，这显然给出如上交换图。

为说明映射 $[X,G_k]\to\Vect_k(X)$ 的满性，设 $p: E\to X$ 是 $k$ 维向量丛，$\{U_i\}$ 是 $X$ 的开覆盖，使 $E$ 在每个 $U_i$ 上平凡。则由 $X$ 的仿紧性，存在 $X$ 的可数开覆盖 $\{U_a\}$，使 $E$ 在每个 $U_a$ 上仍平凡，并且存在单位分解 $\{\phi_a\}$，其中 $\phi_a$ 支撑在 $U_a$ 中。设 $g_a: p^{-1}(U_a)\to\C^k$ 是平凡化 $p^{-1}(U_a)\to U_a\times\C^k$ 与到 $\C^k$ 的投影的复合。映射 $(\phi_a p)g_a$（$v\mapsto\phi_a(p(v))g_a(v)$）延拓为映射 $E\to\C^k$，它在 $p^{-1}(U_a)$ 之外为零。在 $X$ 的每点附近只有有限多个 $\phi_a$ 非零，且至少有一个 $\phi_a$ 非零，于是这些延拓后的 $(\phi_a p)g_a$ 构成映射 $g: E\to(\C^k)^\infty=\C^\infty$ 的坐标，而该映射在每条纤维上是线性单射。

为说明单性，若有两个映射 $f_0,f_1: X\to G_k$ 给出的同构 $E\simeq f_0^*(E_k)$ 与 $E\simeq f_1^*(E_k)$，则如证明的第一部分那样，它们给出在纤维上为线性单射的映射 $g_0,g_1: E\to\C^\infty$。于是断言是：$g_0$ 与 $g_1$ 可通过在纤维上仍为线性单射的映射 $g_t$ 同伦。若如此，则 $f_0$ 与 $f_1$ 通过 $f_t(x)=g_t(p^{-1}(x))$ 同伦。

为证明这一断言，第一步是如下构造同伦 $g_t$：把 $g_0$ 与同伦 $L_t: \C^\infty\to\C^\infty$ 复合，后者定义为
\[ L_t(x_1,x_2,\dots)=(1-t)(x_1,x_2,\dots)+t(x_1,0,x_2,\dots). \]
对每个 $t$ 这都是线性映射，其核容易算出为 $0$，故 $L_t$ 是单射。把 $L_t$ 与 $g_0$ 复合，便把 $g_0$ 的像移入奇数号坐标。类似地我们可以把 $g_1$ 同伦到偶数号坐标。我们仍把新的 $g$ 记作 $g_0$ 与 $g_1$，并令 $g_t=(1-t)g_0+tg_1$。它是线性的，且对一切 $t$ 在纤维上都是单射，因为 $g_0$ 与 $g_1$ 如此。
\end{proof}

% [Tutor Note: p.26 交换图左下竖直箭头是丛投影 $p: E\to X$；而顶行最右的 $\pi: E_k\to\C^\infty$ 由 $\pi(l,v)=v$ 定义（$E_k=\{(l,v): v\in l\}$），二者不要混淆。图中底部两个 $X$ 之间的 $id$ 说明拉回丛 $f^*(E_k)$ 的底空间仍是 $X$，而第二个 $X\to G_k$ 正是分类映射 $f$。]

% [原文 p.27]
这个定理表明，某个固定拓扑空间上的 $k$ 维向量丛由映入 $G_k$ 的映射的同伦类分类；正因如此，$G_k$ 被称为 $k$ 维向量丛的\emph{分类空间}（classifying space），而 $E_k\to G_k$ 被称为\emph{万有向量丛}（universal vector bundle）。然而实情是，这个定理在枚举给定空间上所有不同的（固定维数的）向量丛方面用处有限，因为 $[X,G_k]$ 的显式计算通常超出技术可达的范围。它的重要性更多在于其理论含义，因为它（除其他作用外）可以把一般命题的证明归约到万有丛这一特殊情形。

上述构造也可用于实情形；还有一个关于定向实向量丛的版本（见 [9]）。

我们想通过描述球面上复向量丛的一个构造来结束本节，即使用所谓的\emph{粘合函数}（clutching functions）。这是一个技术要点，下一章我们会用到。

我们想构造 $k$ 维球面 $S^k$ 上的（实或复，这里取复情形）向量丛 $E\to S^k$。把球面写成上半球 $D^k_+$ 与下半球 $D^k_-$ 的并，且 $D^k_+\cap D^k_-=S^{k-1}$。给定映射 $f: S^{k-1}\to GL_n(\C)$，设 $E_f$ 为不交并 $D^k_+\times\C^n\cup D^k_-\times\C^n$ 把 $(x,v)\in\partial D^k_-\times\C^n$ 与 $(x,f(x)(v))\in\partial D^k_+\times\C^n$ 等同所得的商空间。于是有自然投影 $E_f\to S^k$，它是一个 $n$ 维向量丛。上面用到的映射 $f$ 称为\emph{粘合函数}（因为它做的事情本质上正是车辆中离合器的功能）。事实上可以证明映射 $\Phi: f\mapsto E_f$ 给出双射
\[ \Phi: [S^{k-1},GL_n(\C)]\to\Vect_n(S^k) \]
（这个结论的证明可在 [9] 中找到；事实上它是上面定理 1 的一个有趣的特殊情形）。这个双射对实情形并不完全成立，因为 $GL_n(\R)$ 不是道路连通的。不过它对定向实向量丛成立（详见 [9]）。

% [Tutor Note: p.27 粘合函数双射 $\Phi$ 是定理 1（同伦不变性）的一个特例：把 $S^k$ 沿赤道 $S^{k-1}$ 分成两个闭半球，过渡函数 $f$ 的同伦类决定丛的同构类，因为 $D^k_\pm\times\C^n$ 都平凡，全部信息集中在赤道处的过渡函数上。实情形失效的原因是 $GL_n(\R)$ 有两个道路连通分支（$\det>0$ 与 $\det<0$），故须加定向假设才能得到双射。]

